\section{Changes in version 4}\label{app:changes}

Version~3 of this paper was withdrawn on 27 September 2026, together with the audit notes, because of its framing. Version~4 keeps its theorems, proofs and citations, corrects its framing, and adds new results. The changes are the following.

\paragraph{Framing.}
\begin{itemize}
\item The title, abstract and introduction are rewritten. The introduction states the main results as Theorems~I--V, explains what they mean, describes the bridges, credits the contributors, and sets out the four statuses and how the status labels map onto them.
\item Section~\ref{sec:notes} is retitled \emph{The author's notes, April--June 2026}, and the notes are described in the same way throughout.
\item Quotation marks around the terms of the workbench and the notes are replaced by italics.
\item In the map of the archive (§\ref{ssec:archive}) the category \emph{side} is renamed \emph{structural} and defined through the archive's own description of these sections as an interface still to be supplied by a typed map.
\item Statements that structures which do not take primes as input cannot bear on solvability are replaced by what is established about the typed maps that would connect them (Sections~\ref{sec:notes} and~\ref{sec:overlaps}). The negative results of Section~\ref{sec:negative} are stated with their scope, and the table's closing sentence now says which mechanisms are sharp.
\item The drafts that the author's latest record supersedes are treated only for the results they contain that were verified: the paragraph on one such draft, its row in the catalogue, and the row of Section~\ref{sec:negative} on the terminal core as the residual locus are removed. The terminal core is now described by what Theorem~\ref{thm:termcore} locates (Remark~\ref{rem:termcore}).
\item Statements that rest only on a reading pass over the notes are marked \emph{(a reading pass)} and count as not verified here.
\item Section~\ref{sec:overlaps} is rebuilt as a section on the bridges to other areas, with what is established, what is not verified here, and my assessment for each; the table of edges to the author's other workbenches is its last subsection.
\end{itemize}

\paragraph{Corrections of content.}
\begin{itemize}
\item \emph{The minimal model $M(2,7)$.} Version~3 said that the note's identification of the mock theta shadow with the characters of $M(2,7)$ fails, that the note's modified modular pair is not a representation of $SL_2(\ZZ)$, and that no relabelling matches both $S$ and $T$. Re-deriving the statement shows that the identification holds after conjugation and the twist by the character of $\eta^3$ (Theorem~\ref{thm:m27}). What fails is the note's matching of labels through its cubic phase transport: no signed permutation and scalar realise it, and the pair it produces does not satisfy $(ST)^3=S^2$; version~3's statements hold in that scope only.
\item \emph{The Lorentzian note.} Version~3 listed the note's section on modular flow as an exception to the correctness of the notes' structural links, stating that the modular flow is a rotation and not a Lorentz boost. That statement stands (Proposition~\ref{prop:flowboost}(a)). What changes is the verdict on the note: its proposition, that the congruence action of $A_\rho$ is a Lorentz boost, is correct, and Proposition~\ref{prop:flowboost} gives the exact relation between the two.
\item \emph{The sieve.} Version~3 said that a sieve bound of the form $x/(\log x)^\kappa$ follows from the quadratic-residue conditions and that the conditions do not approach a proof. The first statement was not proved there and is removed; Proposition~\ref{prop:densityzero} proves density zero, and Question~\ref{q:qrsurvivors} states what a negative answer would prove.
\item \emph{The shell count as a class function.} Version~3 said that $A_{11}(p)$ takes the values $0,4,6,8,10$ on primes $p\equiv1\pmod{840}$. The count meant is $A_{11}(pa)$ with $a=(p+11)/4$, the number of ordered certificates of the shell; $A_{11}(p)$ itself takes only the values $0$ and $2$ at primes. Section~\ref{ssec:noteserrors} now states it so, with the primes at which each value occurs.
\item \emph{The further certificates of the referee of version~3.} They are now stated with their proofs, and the $23$ primes below $10^7$ are described as the primes $p\equiv1\pmod{24}$ with no certificate of the kinds considered, rather than as survivors of conditions that version~3 did not state.
\item \emph{The Busy-Beaver notes.} Version~3's statement that the residues modulo $107$ are at chance level is replaced by the description of the test and its scope (Section~\ref{ssec:structfacts}), and its statement that the only link with the Busy-Beaver function is the classical one by the exact $\Pi_1$ bridge (Section~\ref{sec:overlaps}).
\end{itemize}

\paragraph{New results.}
Proposition~\ref{prop:densityzero} (density zero of the survivors of the row $p+1$); Theorem~\ref{thm:m27} (the order-$7$ shadow carries the $M(2,7)$ representation); Proposition~\ref{prop:fibonacci} (the order-$5$ block carries the Fibonacci data); Propositions~\ref{prop:dossier} and~\ref{prop:torsor}, which state as propositions, with proofs, two points that version~3 listed as errors; Proposition~\ref{prop:flowboost} (modular flow and boost); and Theorem~\ref{thm:orbits} with Corollary~\ref{cor:orbits} (the orbits of the $(3,4,\infty)$ covers at every level), which answer version~3's question on the covers at even levels. The coefficients of the note's modular differential equation are checked in Section~\ref{ssec:mock}. Section~\ref{sec:overlaps} also states, from the companion note \cite{ApollonianNote}, how rotations and boosts meet the Apollonian group, and, from a reading pass, that the note's odd Ogg phase decomposition holds with the map of Proposition~\ref{prop:torsor}.

\paragraph{Questions.}
Section~\ref{sec:open} is reordered. Version~3's question on the covers at even levels is answered and replaced by the question on their genera; Questions~\ref{q:brieskorn}--\ref{q:pi1} are new.

Apart from the corrections listed above, no statement of version~3 changed its mathematical content.
